
\chapter{A quantum examination}

\textbf{Purpose:} Ask whether a quantum route to a SHA-256 preimage beats running the function, and
answer it with measurements, and not with an appeal to hardware that does not exist.
\textbf{Scope:} \texttt{nonce\_period.py}, \texttt{shor\_trace.py}, \texttt{qubit\_trace.py},
\texttt{exact\_qubits.py}, \texttt{hhl\_nested.py}, \texttt{src/engine/cuda\_miner.cu}.

Every figure below is measured in this tree, derived in one line from something measured here, or
labeled \emph{quoted}. There are two quoted figures, both from the published literature on quantum
preimage attacks, and the text names each where it uses it.

\section{What the function is, computationally}

SHA-256 is arithmetic in the ring modulo \(2^{32}\). That is the definition and not a tuning
choice: widening the word computes a different function and fails silently.

Its per-round work is three operations. Two are three-input boolean functions, \(\mathrm{Ch}\) and
\(\mathrm{Maj}\). The third is rotation, since the \(\Sigma\) and \(\sigma\) mixers are each three
rotates combined by exclusive or. Rotation is the torsion on the state and it is the dominant cost.
Counted from the compiled machine code of the miner in this tree:

\begin{center}
\begin{tabular}{lrrl}
\toprule
opcode & count & share & what it is \\
\midrule
SHF       & 5394 & 29.90\% & funnel shift, the rotation \\
LOP3      & 3625 & 20.09\% & three-input boolean: Ch and Maj are one instruction each \\
IMAD      & 2126 & 11.78\% & \\
IADD3     & 1996 & 11.06\% & three-way add, fused \\
LDL / STL & 0    & 0\%     & no register spills anywhere \\
\bottomrule
\end{tabular}
\end{center}

Occupancy of the mining kernel is \textbf{100\%} at 38 registers, 48 of 48 warps resident. The
classical implementation is at the integer-pipe roofline and its rotations are 30\% of everything,
which is worth establishing before asking whether a different kind of machine helps.

Rotation is also where the classical optimizations stop. SHA-256's rotate distances are 2/13/22,
6/11/25, 7/18 and 17/19, and \textbf{none is byte-aligned}. No rotate therefore collapses into a single
byte-permute.

\section{Shor is the wrong tool, and the reason is measurable}

Shor's algorithm finds the \emph{period} of a function. It prepares a superposition over the domain,
applies the function, and lets the Fourier transform peak on the period.
\texttt{shor\_trace.py} shows exactly that on \(N = 15\) with base 7: peaks at 0, 4, 8, 12 because
\(7^x \bmod 15\) has period 4, and two controlled multiplies that move nothing at all because
\(7^4 \bmod 15\) closes the orbit. The factors fall out by greatest common divisor,
\(\gcd(48,15) = 3\) and \(\gcd(50,15) = 5\).

That run \textbf{factors 15}. It has never touched a nonce, and no run in this tree has.

So the question deciding Shor's relevance is whether the nonce to hash map has a period.
\texttt{nonce\_period.py} measures it on a chain-verified header with an exact Poisson tail and a
family-wise threshold: \textbf{null across 75 tests, detection limit \(1.83 \times 10^{-2}\)}.

Two walls follow and the second does not care about the first.

\subsection{Size}

Shor needs about \(2n+3\) qubits for an \(n\)-bit domain:

\begin{center}
\begin{tabular}{rrrr}
\toprule
domain bits & qubits & amplitudes & bytes \\
\midrule
16 & 35 & \(10^{11}\) & \(10^{12}\) \\
\textbf{32} & \textbf{67} & \(\mathbf{10^{20}}\) & \(\mathbf{10^{21}}\) \\
\bottomrule
\end{tabular}
\end{center}

A 32-bit nonce space wants 67 qubits as a full state vector, which is \(10^{21}\) bytes. The largest
full state-vector simulations performed anywhere are near 46 qubits.

\subsection{Structure}

Hand a 67-qubit machine an aperiodic function and it returns a flat distribution, correctly and
uselessly. The resource count is not the obstruction. The absent period is.

\section{What a qubit costs, and where the exponential lives}

\begin{center}
\begin{tabular}{lrrr}
\toprule
representation & bytes a qubit & 68 qubits & how many in 46 GB \\
\midrule
product state, \(4 \times\) float64 & 32 & 2176 B & \(2.29 \times 10^{7}\) \\
graph state, 4 terms and adjacency & \(\approx 36\) & 2461 B & \(2.03 \times 10^{7}\) \\
full state vector & not linear & \(4.72 \times 10^{21}\) B & \textbf{0} \\
\bottomrule
\end{tabular}
\end{center}

Sixty-eight qubits is 2.46 kB in the graph form and \(4.72 \times 10^{21}\) bytes as a state vector.
The same qubit count, eighteen orders of magnitude apart. \textbf{The exponential lives in the
entanglement and never in the qubit count.}

The cheap representations are cheap for a reason with weight here: they hold stabilizer states, and
stabilizer states are classically simulable. \texttt{exact\_qubits.py} measures the same wall from
the other side. The Clifford gate set is exact at \(0.000 \times 10^{0}\) at every width and every
depth, because those gates are swaps and sign flips with no arithmetic in them. Nothing that exact
is doing work.

Sixty-four bits a term is the floor, measured and not preferred. At \(2 \times\)
float32 the state error is \(1.2157 \times 10^{-7}\); at \(2 \times\) float16 it is
\(1.0162 \times 10^{-3}\), which sits above the detection limits this tree's own instruments use. A
representation cheaper than the measurements it serves is not a saving.

\section{Verdict}

\textbf{The proof of work is structurally resistant, and the argument contains no hardware
estimate.} Grover gives \(\sqrt{N}\) at best, the BBBV bound says nothing does better on a
structureless oracle, four instruments here found no structure to exploit, the oracle is a coherent
double SHA-256, and parallel Grover buys \(\sqrt{m}\) where classical buys \(m\). Adding hardware
helps the defenders faster than the attacker.

\textbf{The signatures are a different matter.} Bitcoin spends are secured by ECDSA on secp256k1,
and Shor breaks discrete logarithms outright given a large enough fault-tolerant machine. That is
the same Shor as \texttt{shor\_trace.py}, applied to a problem that \emph{does} carry the periodic
structure the nonce map lacks. The deterrent there is that the machine does not exist yet, a
time sink, and not a lock.

So the mining is safe and the coins are not, and those two facts get called quantum resistance
interchangeably when only one of them holds.

The lever on the mining side is unchanged by any of this. The card exhausts a 32-bit nonce space in
\textbf{1.82 seconds} at \(2.359 \times 10^{9}\) hashes a second. More trials is the only thing that
has ever moved the number.
